\section{A threshold and a critical exponent have different premises}
A smooth trajectory of density matrices can cross the purity cut without a change of stability. In Section~\ref{sec:geometry}, choose $t$ increasing smoothly through $1/\sqrt6$ in the family $\rho_t$. The majority test $P>2/7$ changes truth value while $\rho_t$ and its eigenvalues remain smooth. The equation for $P$ contains no singularity there.

There is a separate conditional calculation for a scalar mean-field potential \citep{griffiths1970,lawrie1984}. Supply an order parameter $m$, control $t$, field $h$, and coefficients $v>0$ and $u$ in
\[
 F(m;t,h)=\tfrac12 t m^2+\tfrac14 u m^4+\tfrac16 v m^6-hm .
\]
Under the tuning $u=0$, the stationary equation is $tm+vm^5=h$. At $h=0$ and $t<0$, the nonzero minimizing amplitudes obey $|m|=(-t/v)^{1/4}$, so $\beta=1/4$. At $t=0$ one has $|h|=v|m|^5$, giving $\delta=5$. Above the transition, the zero-field susceptibility is $\partial m/\partial h=1/t$, giving $\gamma=1$. Below it the minimized free energy is $-|t|^{3/2}/(3\sqrt v)$, whose second derivative scales as $|t|^{-1/2}$ and gives $\alpha=1/2$. For $u>0$ the leading nonzero amplitude near zero instead scales as $(-t/u)^{1/2}$. The exponents therefore depend on the supplied symmetry, tuning and mean-field approximation. Relating $m$ to an independently measured quantity, or $t$ to purity in an actual device, requires an additional model and evidence. No physiological or clinical exponent is reported here.

\paragraph{Reproducibility.}
The linked proof sources give executable statements for the operational quotient, transfer, address and rebit tests. The scalar window and controls are also documented in the \href{https://holon.sh/docs/consciousness/hierarchy/interiority-hierarchy}{model hierarchy} and \href{https://holon.sh/docs/reference/mathematical-kernel}{typed mathematical kernel}. The displayed matrix and mean-field bounds are derived above. The proposed binocular-rivalry test requires an independently registered protocol and participant data before it can yield an empirical conclusion.
